\documentclass[10pt,a4paper]{amsart}
\usepackage[margin=19mm]{geometry}
\usepackage{amsmath,amssymb,mathtools}
\usepackage[scr=rsfs,cal=euler]{mathalfa}
\usepackage{xfrac}
\usepackage[unicode,hidelinks,bookmarks=false]{hyperref}
\newcommand{\ie}{i.e.\ }
\newcommand{\cf}{cf.\ }
\newcommand{\resp}{respectively\ }
\newcommand{\Subsection}{\subsection}
\providecommand{\nobreakdash}{\nobreak\hbox{-}}
\setlength{\emergencystretch}{2em}
\multlinegap0pt
\usepackage{amssymb}
\usepackage{csquotes}
\usepackage{mathtools}
\usepackage{braket}
\usepackage{dsfont}
\usepackage[shortlabels]{enumitem}
\usepackage{caption}
\usepackage{subfig}
\usepackage[all]{xy}
\usepackage{xfrac}
\usepackage{bbm}
\usepackage{esint}
\usepackage{tikz}
\newtheorem{lemma}{Lemma}
\newcommand{\N}{\mathbb{N}}
\newcommand{\PP}{\mathcal{P}}
\newcommand{\R}{\mathbb{R}}
\title{Nested superposition principle for random measures and the geometry of the Wasserstein on Wasserstein space}
\author{Alessandro Pinzi, Giuseppe Savar\'e}
\date{Source: arXiv:2510.07523v2 (CC BY 4.0)}
\begin{document}
\maketitle

\noindent\emph{Source and licence.} Nested superposition principle for random measures and the geometry of the Wasserstein on Wasserstein space. Version arXiv:2510.07523v2 (CC BY 4.0), distributed by arXiv under the Creative Commons Attribution 4.0 International licence, http://creativecommons.org/licenses/by/4.0/, as recorded in the arXiv licence metadata for this exact version and shown on the abstract page https://arxiv.org/abs/2510.07523v2. Full licence notice: \texttt{licenses/arxiv-2510.07523-CC-BY-4.0.txt}.\par\smallskip

\noindent\textit{Editorial scope.} Counted unit: the article's lemma infty (source0.tex lines 3591--3601). The article's complete original proof of it is delivered with it (source0.tex lines 3603--3655, one \texttt{proof} environment of 53 lines). No mathematics is written, repaired or re-derived: the statement and the proof are the article's own lines, in the article's own notation, and no retained line is substituted or re-flowed.  No further source-local context is needed: every cross-reference of every retained line resolves inside the two delivered spans.  Nothing was omitted from any retained span, so the package carries no omission record.  No retained line contains a citation, so the wrapper carries no bibliography and no bibliography entry.  No retained line contains a figure environment, an image-inclusion command or drawing code, so no image is delivered and the package declares no asset dependency.  Editorial formatting only, in addition: an AMS \texttt{amsart} preamble that restates the article's own declarations and macros used by the retained lines, namely the environments \texttt{lemma} on separate counters, together with the standard packages those lines need.  The retained environments print in this document as Lemma 1 in order of appearance, whereas the article prints the same items with the numbering of its own full document.  The complete native source of the article as distributed in the arXiv e-print for this exact version is bundled unchanged as \texttt{sources/186680/source0.tex}.
\begin{lemma}\label{lemma: dual cyl R infty}
    For any $\mathtt{m}_0,\mathtt{m}_1\in \PP(\R^\infty)$ it holds
    \begin{equation}\label{dual formula R^infty}
    \begin{aligned}
        W_{1,D_\infty} &(\mathtt{m}_0,\mathtt{m}_1) =  \sup \bigg\{ \int_{\R^\infty} F d(\mathtt{m}_0-\mathtt{m}_1) \ : \ F(x)\in\operatorname{Cyl}_c^1(\R^\infty), \ F \ 1\text{-}\operatorname{Lip} \text{ w.r.t. } D_\infty\bigg\}
        \\
        = & \sup \bigg\{ \int_{\R^\infty}\hspace{-0.2cm} \Psi\circ \pi_k d(\mathtt{m}_0-\mathtt{m}_1)  :  k\in \N, \, \Psi \in C_c^1(\R^k), \, \|\Psi\|_\infty\leq 1/2,\, \bigg\|\sum_{i=1}^k|\partial_i\Psi|\bigg\|_\infty \leq 1\bigg\}
        .
    \end{aligned}
    \end{equation}
\end{lemma}

\begin{proof}
We prove the first inequality, while the second one will be a byproduct of our argument.
    The $\geq$ inequality is trivial, so we focus on the other one. Define the projection functions
    \[\pi^n:\R^\infty \to \R^\infty, \quad \pi^n(\mathrm{x}) = (x_1,\dots,x_n,0,0,\dots)\]
    \[ \hspace{-1,3cm}p^n:\R^\infty\to\R^n, \quad p^n(\mathrm{x})= (x_1,\dots,x_n).\]
    Notice that $W_{1,D_\infty}(\pi^n_\sharp \mathtt{m}_0,\pi^n_\sharp \mathtt{m}_1)\leq W_{1,D_\infty}(\mathtt{m}_0,\mathtt{m}_1)$. Moreover, for all $\mathtt{m}\in \PP(\R^\infty)$, we have $\pi^n_\sharp \mathtt{m}\to\mathtt{m}$ weakly in duality with $C_b(\R^\infty,\tau_w)$. Indeed, for any $\mathrm{x}\in \R^\infty$, $\pi^n(\mathrm{x})\overset{\tau}{\to}\mathrm{x}$, so for any $F\in C_b(\R^\infty,\tau_w)$, by the dominated convergence theorem, it holds
    \[\int F(\mathrm{x}) d(\pi^n_\sharp \mathtt{m})(\mathrm{x}) = \int F(\pi^n(\mathrm{x})) d\mathtt{m}(\mathrm{x}) \to \int F(\mathrm{x}) d\mathtt{m}(x).\]
    Moreover, $D_\infty$ is $\tau_w \otimes \tau_w$ l.s.c.: indeed consider $\mathrm{x}^{(n)},\mathrm{y}^{(n)}\in \R^\infty$ respectively, converging component-wise to $\mathrm{x}$ and $\mathrm{y}$, then 
    \begin{align*}
        \liminf_{n\to+\infty} D_\infty(\mathrm{x}^{(n)},\mathrm{y}^{(n)}) 
        = &
        \liminf_{n\to+\infty} \sup_{j\in \N} |x_j^{(n)} - y_j^{(n)}|\wedge 1 
        \\
        \geq &
        \liminf_{n\to+\infty} \sup_{j\leq m} |x_j^{(n)} - y_j^{(n)}|\wedge 1 = \sup_{j\leq m} |x_j - y_j|\wedge 1,
    \end{align*}
    for any $m\in \N$, and for its arbitrariness we conclude. 
    \\
    Consider a sequence of $W_{1,D_\infty}$-optimal plans $\Pi_n \in \Gamma_0(\pi^n_\sharp \mathtt{m}_0,\pi^n_\sharp \mathtt{m}_1)$. By the convergence of the marginals and the fact that $(\R^\infty,\tau_w)$ is a Polish space, we have that the set of measures $\{\Pi_n\}$ is tight, which implies that 
    \[\forall (n_k) \ \exists (n_{k_j}) \ : \ \Pi_{n_{k_j}} \to \Pi\in\Gamma(\mathtt{m}_0,\mathtt{m}_1),\]
    where the convergence is weakly in duality with $C_b(\R^\infty\times\R^\infty, \tau\otimes\tau)$. Then, by $\tau_w\otimes\tau_w$-l.s.c. of $D_\infty$ we have 
    \begin{align*}
        \liminf_{j\to+\infty} W_{1,D_\infty}(\pi^{n_{k_j}}_\sharp \mathtt{m}_0,\pi^{n_{k_j}}_\sharp \mathtt{m}_1) 
        = &
        \liminf_{j\to+\infty} \int D_\infty(\mathrm{x},\mathrm{y}) d\Pi_{n_{k_j}}(\mathrm{x},\mathrm{y}) 
        \\
        \geq & 
        \int D_\infty(\mathrm{x},\mathrm{y}) d\Pi(\mathrm{x},\mathrm{y}) \geq W_{1,D_\infty}(\mathtt{m}_0,\mathtt{m}_1),
    \end{align*}
    and by arbitrariness of $(n_k)$ we conclude that $\liminf_{n\to+\infty} W_{1,D_\infty}(\pi^{n}_\sharp \mathtt{m}_0,\pi^{n}_\sharp \mathtt{m}_1) 
        \geq W_{1,D_\infty}(\mathtt{m}_0,\mathtt{m}_1) 
        $. Together with what was proved before, we have that 
        \begin{equation}
            \lim_{n\to+\infty} W_{1,D_\infty}(\pi^{n}_\sharp \mathtt{m}_0,\pi^{n}_\sharp \mathtt{m}_1) 
        = W_{1,D_\infty}(\mathtt{m}_0,\mathtt{m}_1) .
        \end{equation}
        Now, for any $x,y\in \R^n$ define $D_n(x,y) := \sup_{j\leq n} |x_j-y_j|\wedge 1$ and notice that 
        \begin{align*}
            W_{1,D_\infty}(\mathtt{m}_0,\mathtt{m}_1) 
            = & 
            \sup_{n\in \N} W_{1,D_\infty}(\pi^n_\sharp \mathtt{m}_0,\pi^n_\sharp \mathtt{m}_1)
            = 
            \sup_{n\in \N} W_{1,D_n}(\pi^n_\sharp \mathtt{m}_0,\pi^n_\sharp \mathtt{m}_1)
            \\
            = &
            \sup_{n\in \N} \ \sup_{\{\Psi\in C_c^1(\R^n): \Psi\  D_n \ 1-\operatorname{Lip}\}} \left\{ \int_{\R^n} \Psi(x) d\pi^n_\sharp \mathtt{m}_0(x) - \int_{\R^n} \Psi(y) d\pi^n_\sharp \mathtt{m}_1(y) \right\}
            \\
            = & \sup_{n, \Psi \text{ as above}} \int \Psi(\pi^n(\mathrm{x})) d(\mathtt{m}_0 - \mathtt{m}_1)(\mathrm{x}).
        \end{align*}
        The proof is then concluded by the fact that $\R^\infty \ni x \mapsto \Psi(\pi^n(x))$ is the general expression for a cylinder function from $\R^\infty\to\R$, and it is $D_\infty$ $1$-Lipschitz if and only if $\R^n \ni x \to \Psi(x)$ is $D_n$ $1$-Lipschitz. This is equivalent to ask that 
        \[\left\|\sum_{i=1}^n |\partial_i\Psi|\right\|_\infty  \leq 1 \quad \text{ and }\quad \operatorname{osc}\Psi:= \max \Psi - \min \Psi \leq 1.\]
        Moreover, up to considering a translation, we can substitute the condition on oscillation with $\|\Psi\|_\infty\leq 1/2$, which concludes the proof.
\end{proof}

\end{document}
